\lesson{7}{Mar 10 2022 Thu (16:13:36)}{Surprise}{Unit 3}

\subsubsection{Twist Ending Techniques}
\label{sub_sub_sec:twist_ending_techniques}

In a typical story, the plot or the story line often follows a predictable
chronological order. However, in suspenseful stories, the order of events may
be changed to build suspense and create questions for audiences.

Suspenseful stories still contain an exposition, rising action, climax, falling
action, and a resolution. However, authors use various techniques to change the
sequence of events to surprise, challenge, and engage audiences. The exposition
might start with the action with what happens last in the story, and the
resolution might reveal the very first action that  set the whole frightening
plot in motion.

Here are some twist ending techniques:

\begin{definition}[Discovery]
  \label{def:discovery}

  The protagonist's sudden recognition of his or her own identity or nature or
  another character's identity or nature.

  The audience learns the true identity or nature the moment the character
  discovers it.
\end{definition}

\begin{definition}[Flashback]
  \label{def:flashback}

  A sudden switch to a past event that provides a missing piece that solves the
  mystery.

  The audience learns previously unknown information, which reveals the reason
  for something previously unexplained or puts a character or situation in a
  different light.
\end{definition}

\begin{definition}[Foreshadowing]
  \label{def:foreshadowing}

  A character or plot element that is introduced early in the story; it may not be
  referenced until later, or it may be something that appears throughout the story
  or film but whose significance is not realized until later.

  Audiences might suspect an individual or wonder about the reappearance of a
  particular object or sound, but they do not realize the significance until the
  end.
\end{definition}

\begin{definition}[In medias res]
  \label{def:in_medias_res}

  A technique in which the story begins in the middle rather than the beginning.
  In medias res is Latin for \textit{"into the middle of things"}.

  Audiences rely on flashbacks to learn about setting, characterization, and
  conflict.
\end{definition}

\begin{definition}[Nonlinear Narration]
  \label{def:nonlinear_narration}

  A method of telling the story out of order.

  The audience has to put the pieces the pieces of the story together. The
  information is withheld strategically to build tension. As the pieces of the
  puzzle come together, events or characters might take on different meanings.
\end{definition}

\begin{definition}[Reversal]
  \label{def:reversal}

  A change of the protagonist's fate that comes about because of a twist in the
  character's circumstances.

  Audiences are surprised by the change in circumstances. For example, we might
  discover that the whole story has been a dream, a simulation, or an act. We
  might also discover that the character has a twin, is actually sick, or is
  really the villain.
\end{definition}

\begin{definition}[Reverse Chronology]
  \label{def:reverse_chronology}

  The revelation of the plot in reverse order, from the end to the beginning.

  Audiences know the ending of the story. The twist ending is actually the
  beginning, which reveals the cause of all of the action.
\end{definition}

\begin{definition}[Unreliable Narrator]
  \label{def:unreliable_narrator}

  The revelation that the narrator has manipulated or manufactured the story. It
  might also be revealed that the narrator is not sane and therefore offers a
  view of the story that is questionable.

  Audiences question what they accepted to be true of the text.
\end{definition}

% subsubsection twist_ending_techniques (end)

\newpage
